
Annotation-free methods for detecting spurious correlations in neural networks rely on clustering model representations to discover spurious groups, yet the conditions under which this detection step succeeds have not been characterized. This work examines whether clustering-based spurious direction recovery is feature-strength conditional. Applying k-means clustering (k=2, n_init=10, random seed 42) to epoch-5 penultimate-layer embeddings of a pretrained ResNet-50 yields Adjusted Mutual Information (AMI) of 0.762 and worst-cluster purity of 0.892 on Waterbirds, but AMI of 0.258 and purity of 0.456 on CelebA — a ratio of approximately 3:1 in AMI under identical experimental conditions. The CelebA purity of 0.456 is 0.016 above the random baseline of 0.440, representing performance that is statistically near-chance. The key factor differentiating the two outcomes is whether the spurious attribute is encoded as a separable dimension in the ImageNet pretraining prior: scene-level background features (Waterbirds) are pre-encoded in the pretrained weights and are cluster-separable at epoch 5, whereas fine-grained hair texture attributes (CelebA) are not represented in the pretrained prior at the required level of separability, and five epochs of empirical risk minimization (ERM) fine-tuning are insufficient to establish this structure. The paper characterizes this conditionality, identifies the mechanism through pretraining alignment, and proposes a linear probe pre-screen as a candidate diagnostic, noting that the downstream Gradient SNR Balancing (GSB) intervention was not evaluated due to the CelebA detection failure blocking the experiment pipeline. A gap is identified in how published clustering-based detection results are reported across benchmark configurations.

